% =========================================================================
\section{Campos escalares, diferencial y gradiente}\label{sec:campos}
% =========================================================================

% -------------------------------------------------------------------------
\subsection{Campos escalares y su diferencial}\label{sec:escalares}
% -------------------------------------------------------------------------
\begin{definicion}[Campo escalar]\label{def:escalar}
Sea $W\subseteq\R^{n}$ una región. Un \emph{campo escalar} es una función
$f:W\to\R$ que asigna a cada punto $\vc{r}=(x,y,z)\in W$ un número
$f(\vc{r})=f(x,y,z)$.
\end{definicion}

Son campos escalares la temperatura $T(\vc{r})$ en el interior de una
habitación, la altura $h(x,y)$ sobre el nivel del mar en un mapa topográfico
(aquí $n=2$) o el potencial electrostático $V(\vc{r})$. En una sola
variable, la diferencial de $f$ recoge cuánto cambia la función cuando la
variable cambia en una cantidad infinitesimal:
\begin{equation}\label{eq:df1d}
  \dif f=\frac{\dif f}{\dif x}\,\dif x
\end{equation}
(\cref{fig:diferencial1d}): es el incremento a lo largo de la \emph{recta
tangente}, de pendiente $\tan\alpha=f'(x)$. Para un campo escalar de tres
variables hay tres direcciones independientes en las que moverse y la
diferencial total es la suma de las tres contribuciones,
\begin{equation}\label{eq:df3d}
  \dif f=\pd{f}{x}\,\dif x+\pd{f}{y}\,\dif y+\pd{f}{z}\,\dif z .
\end{equation}

\begin{figure}[ht]
\centering
\begin{tikzpicture}[>=Stealth, x=1cm, y=1cm]
  % ejes
  \draw[eje] (-0.3,0) -- (7.6,0) node[right, text=black] {$x$};
  \draw[eje] (0,-0.3) -- (0,4.6) node[above, text=black] {$f$};
  % curva
  \draw[azulprof, line width=1.2pt, domain=0.4:7.2, samples=100, smooth]
    plot (\x, {0.55+3.6/(1+exp(-1.35*(\x-3.7)))});
  % punto, tangente y triángulo diferencial
  \def\xo{3.2}
  \pgfmathsetmacro{\yo}{0.55+3.6/(1+exp(-1.35*(\xo-3.7)))}
  \pgfmathsetmacro{\sl}{3.6*1.35*exp(-1.35*(\xo-3.7))/(1+exp(-1.35*(\xo-3.7)))^2}
  \draw[azulacero, dashed, line width=0.7pt, domain=\xo-1.7:\xo+2.4]
    plot (\x, {\yo+\sl*(\x-\xo)});
  \fill[azulprof] (\xo,\yo) circle (2pt);
  \draw[grispizarra, densely dashed] (\xo,0) -- (\xo,\yo) -- (0,\yo);
  \draw[rojoteo, line width=0.9pt] (\xo,\yo) -- node[below, text=rojoteo] {$\dif x$} (\xo+1.3,\yo);
  \draw[rojoteo, line width=0.9pt] (\xo+1.3,\yo) -- node[right, text=rojoteo] {$\dif f$}
        (\xo+1.3,{\yo+1.3*\sl});
  \draw[grispizarra, densely dashed] (\xo+1.3,0) -- (\xo+1.3,\yo);
  \node[below, font=\small] at (\xo,0) {$x$};
  \node[below, font=\small] at (\xo+1.3,0) {$x+\dif x$};
  \draw[grispizarra, thick] (\xo+0.75,\yo) arc (0:{atan(\sl)}:0.75);
  \node[font=\small, text=black!80] at (\xo+1.18,\yo+0.25) {$\alpha$};
  \node[azulacero, font=\small, anchor=north east] at (\xo-0.55,{\yo-0.55*\sl-0.1}) {$\tan\alpha=f'(x)$};
  \node[azulprof, font=\small, anchor=south] at (6.7,4.3) {$f(x)$};
\end{tikzpicture}
\caption{Diferencial de una función de una variable: $\dif f=f'(x)\,\dif x$
es el incremento a lo largo de la tangente, y coincide con el incremento real
de $f$ salvo infinitésimos de orden superior.}
\label{fig:diferencial1d}
\end{figure}

Si escribimos el desplazamiento infinitesimal
$\dl=\dif x\,\ux+\dif y\,\uy+\dif z\,\uz$ (el \emph{elemento de línea}, un
vector), la \cref{eq:df3d} puede leerse como un producto escalar:
\begin{equation}\label{eq:df-gradiente}
  \dif f=\Bigl(\pd{f}{x}\,\ux+\pd{f}{y}\,\uy+\pd{f}{z}\,\uz\Bigr)\cdot
         \bigl(\dif x\,\ux+\dif y\,\uy+\dif z\,\uz\bigr).
\end{equation}
El primer factor es el protagonista de esta sección.

\begin{definicion}[Gradiente]\label{def:gradiente}
Sea $f:W\subseteq\R^{3}\to\R$ un campo escalar derivable. Su \emph{gradiente}
es el campo vectorial
\[
  \nab f\equiv\grad f=\pd{f}{x}\,\ux+\pd{f}{y}\,\uy+\pd{f}{z}\,\uz,
\]
y la diferencial se escribe, de forma compacta e independiente de la
descripción por componentes,
\begin{equation}\label{eq:df-compacta}
  \boxed{\;\dif f=\nab f\cdot\dl\;}
\end{equation}
\end{definicion}

\begin{observacion}[El gradiente no depende de las coordenadas]\label{obs:grad-intrinseco}
La fórmula de la \cref{def:gradiente} parece depender de los ejes elegidos, pero
no es así: el gradiente es el único vector que satisface, para todo vector
$\vc{v}$, la relación $\nab f(P)\cdot\vc{v}=D_{\vc{v}}f(P)$ con la derivada
direccional que definimos a continuación~\cite{NAVAS}. En
física se acostumbra a escribir los campos conservativos como $\vc{F}=-\nab f$
cuando $f$ es una energía potencial; aquí no cambiamos el signo.
\end{observacion}

% -------------------------------------------------------------------------
\subsection{Derivada direccional y significado geométrico}\label{sec:direccional}
% -------------------------------------------------------------------------
Necesitamos primero una versión de la regla de la cadena: cómo cambia $f$ si
nos movemos a lo largo de una curva.

\begin{proposicion}[Derivada de un campo escalar a lo largo de una curva]\label{prop:cadena}
Sea $f$ un campo escalar de clase $C^{1}$ y $\gamma:[a,b]\to W$ una curva
derivable. Entonces
\begin{equation}\label{eq:cadena}
  (f\circ\gamma)'(t)=\nab f\bigl(\gamma(t)\bigr)\cdot\gamma'(t).
\end{equation}
\end{proposicion}
\begin{proof}
La matriz jacobiana de $f$ es la matriz fila $\bigl(\partial_xf,\partial_yf,\partial_zf\bigr)$,
la de $\gamma$ es la columna $\bigl(x'(t),y'(t),z'(t)\bigr)^{\!\top}$ y la de la
composición $f\circ\gamma$ es la matriz $1\times1$ formada por su derivada
ordinaria. La regla de la cadena multivariable dice que la jacobiana de la
composición es el producto de las dos primeras, es decir,
$(f\circ\gamma)'=\partial_xf\,x'+\partial_yf\,y'+\partial_zf\,z'$, que es
$\nab f\cdot\gamma'$.
\end{proof}

\begin{definicion}[Derivada direccional]\label{def:direccional}
Sea $\hat u$ un vector unitario. La \emph{derivada direccional} de $f$ en el
punto $P$ en la dirección $\hat u$ es el ritmo de variación de $f$ al avanzar
sobre la recta que pasa por $P$ con dirección $\hat u$:
\[
  D_{\hat u}f(P)\defeq\lim_{h\to0}\frac{f(P+h\,\hat u)-f(P)}{h}.
\]
\end{definicion}

\begin{proposicion}\label{prop:direccional}
$D_{\hat u}f(P)=\nab f(P)\cdot\hat u=\norm{\nab f(P)}\cos\epsilon$, donde
$\epsilon$ es el ángulo que forman $\nab f(P)$ y $\hat u$.
\end{proposicion}
\begin{proof}
Aplicamos la \cref{prop:cadena} a la recta $\gamma(t)=P+t\,\hat u$, cuya
velocidad es $\gamma'(t)=\hat u$: $D_{\hat u}f(P)=(f\circ\gamma)'(0)=\nab f(P)\cdot\hat u$.
La segunda igualdad es la definición geométrica del producto escalar con
$\norm{\hat u}=1$.
\end{proof}

\begin{proposicion}[Dirección de máximo crecimiento]\label{prop:maximo}
Sea $\nab f(P)\neq\vc{0}$. Entre todas las direcciones unitarias,
\[
  \max_{\norm{\hat u}=1}D_{\hat u}f(P)=\norm{\nab f(P)}
  \qquad\text{y se alcanza si y solo si}\qquad
  \hat u=\frac{\nab f(P)}{\norm{\nab f(P)}}.
\]
\end{proposicion}
\begin{proof}
Por la desigualdad de Cauchy--Schwarz,
$D_{\hat u}f(P)=\nab f(P)\cdot\hat u\le\norm{\nab f(P)}\,\norm{\hat u}=\norm{\nab f(P)}$,
con igualdad exactamente cuando $\hat u$ es paralelo y del mismo sentido que
$\nab f(P)$, es decir, cuando $\cos\epsilon=\cos0=1$.
\end{proof}

\begin{observacion}[Las dos preguntas que responde el gradiente]\label{obs:dos-preguntas}
La \cref{prop:maximo} dice que en cada punto el gradiente contesta a dos
preguntas a la vez:
\[
  \nab f\;\longrightarrow\;
  \begin{cases}
    \text{\emph{dirección} en la que $f$ crece más deprisa,}\\[2pt]
    \norm{\nab f}=\text{\emph{valor} de esa tasa máxima de crecimiento.}
  \end{cases}
\]
En la dirección opuesta, $-\nab f$, la función decrece lo más rápido posible,
y en las direcciones perpendiculares a $\nab f$ no cambia (a primer orden).
En términos de la \cref{eq:df-compacta}: $\dif f=\norm{\nab f}\,\norm{\dl}\cos\theta$ (con $\theta$ el ángulo entre $\nab f$ y $\dl$),
y por tanto, para un desplazamiento de módulo fijo, $\dif f$ es máximo cuando
$\theta=0$, esto es, cuando $\dl$ es paralelo a $\nab f$ y entonces
$\dif f=\norm{\nab f}\,\dif\ell$.
\end{observacion}

\begin{proposicion}[El gradiente es ortogonal a las superficies de nivel]\label{prop:ortogonal}
Sea $c\in\R$ y $\Sigma_c=\{\vc{r}\in W: f(\vc{r})=c\}$ la superficie de nivel
(\emph{equipotencial}, o curva de nivel si $n=2$). Si $\gamma$ es una curva
contenida en $\Sigma_c$ con $\gamma(t_0)=P$, entonces
$\nab f(P)\cdot\gamma'(t_0)=0$: el gradiente es perpendicular a $\Sigma_c$ en $P$.
\end{proposicion}
\begin{proof}
Como $f(\gamma(t))=c$ para todo $t$, su derivada es nula. Por la
\cref{prop:cadena}, $0=(f\circ\gamma)'(t_0)=\nab f(P)\cdot\gamma'(t_0)$, y
$\gamma'(t_0)$ recorre todos los vectores tangentes a $\Sigma_c$ en $P$.
\end{proof}

Geométricamente (\cref{fig:gradiente-geo,fig:gradiente-superficie}): si $h(x,y)$ es la altura de un
terreno, sus curvas de nivel son las líneas de «misma cota» del mapa
topográfico, y el gradiente apunta cuesta arriba, cortando a esas curvas en
ángulo recto. Es la línea de máxima pendiente: la que seguiría una gota de agua
(en sentido opuesto) al deslizarse por la ladera.

% Figura: interpretación geométrica del gradiente
\begin{figure}[ht]
\centering
\begin{minipage}[t]{0.47\textwidth}\centering
\begin{tikzpicture}[>=Stealth, baseline=(current bounding box.north)]
  \coordinate (P) at (0,0);
  \draw[vect, line width=1.4pt, acento] (P) -- (4.0,1.6) node[above right, text=acento] {$\nab f$};
  \draw[vect, line width=1.2pt, azulprof] (P) -- (2.2,-0.9) node[right, text=azulprof] {$\dl$};
  % ángulo
  \draw[grispizarra, line width=0.5pt] (1.1,{1.1*1.6/4.0}) arc (21.8:-22.2:1.1);
  \node[font=\small] at (1.45,0.0) {$\theta$};
  % proyección de dl sobre grad f
  \pgfmathsetmacro{\pr}{(2.2*4.0-0.9*1.6)/(4.0*4.0+1.6*1.6)}
  \draw[densely dashed, grispizarra] (2.2,-0.9) -- ({\pr*4.0},{\pr*1.6});
  \draw[line width=2.4pt, azulclaro] (P) -- ({\pr*4.0},{\pr*1.6});
  \node[font=\footnotesize, text=azulacero, rotate=21.8, anchor=south] at ({\pr*2.0},{\pr*0.8}) {$\dl\cos\theta$};
  \draw[grispizarra, line width=0.5pt] ({\pr*4.0-0.18*0.928},{\pr*1.6-0.18*0.371})
     -- ++(0.371*0.18,-0.928*0.18) -- ++(0.928*0.18,0.371*0.18);
  \fill (P) circle (1.6pt);
\end{tikzpicture}\\[4pt]{\small\scshape (a)}\; $\dif f=\norm{\nab f}\,\norm{\dl}\cos\theta$
\end{minipage}\hfill
\begin{minipage}[t]{0.47\textwidth}\centering
\begin{tikzpicture}[>=Stealth, baseline=(current bounding box.north), scale=0.9, transform shape]
  \foreach \c/\pc in {1/100,2/65,3/35}{%
    \draw[azulprof!\pc!azulclaro, line width=0.9pt] (0,0) ellipse ({1.0*\c} and {0.5*\c});}
  \fill[azulprof] (0,0) circle (1.4pt);
  
  % gradientes en puntos de las elipses (apuntan hacia el centro)
  \foreach \c/\t in {2/60,2/200,2/320}{%
    \pgfmathsetmacro{\px}{1.0*\c*cos(\t)}\pgfmathsetmacro{\py}{0.5*\c*sin(\t)}%
    \pgfmathsetmacro{\gx}{-\px/1.0}\pgfmathsetmacro{\gy}{-4*\py/1.0}% -grad(x^2/1+4y^2)/... (dir. al centro)
    \pgfmathsetmacro{\gn}{sqrt(\gx*\gx+\gy*\gy)}%
    \pgfmathsetmacro{\ux}{0.62*\gx/\gn}\pgfmathsetmacro{\uy}{0.62*\gy/\gn}%
    \draw[vect, acento, line width=1pt] (\px,\py) -- ++(\ux,\uy);
    \fill (\px,\py) circle (1pt);}
  \node[font=\footnotesize, text=black!80, anchor=west] at (3.2,1.3) {curvas de nivel};
  
  \node[font=\footnotesize, text=azulprof, anchor=west] at (3.2,-1.3) {$h_3<h_2<h_1$};
\end{tikzpicture}\\[4pt]{\small\scshape (b)}\; $\nab h\perp\{h=\text{cte}\}$
\end{minipage}
\caption{Interpretación geométrica del gradiente. (a)~El incremento
$\dif f$ es el producto de $\norm{\nab f}$ por la proyección del desplazamiento
sobre la dirección del gradiente: es máximo si $\dl\parallel\nab f$ y nulo si
$\dl\perp\nab f$. (b)~Mapa topográfico de una colina: el gradiente de la altura
$h$ es perpendicular a las curvas de nivel y apunta hacia la cumbre (máxima
pendiente).}
\label{fig:gradiente-geo}
\end{figure}


% Figura: el gradiente sobre la superficie z=f(x,y): dirección de máximo ascenso
\begin{figure}[ht]
\centering
\begin{tikzpicture}
  \begin{axis}[superficie, width=0.92\textwidth, height=0.66\textwidth,
      xlabel={$x$}, ylabel={$y$}, zlabel={$z=f(x,y)$},
      xmin=-2, xmax=2, ymin=-2, ymax=2, zmin=0, zmax=1.15,
      xtick={-2,-1,0,1,2}, ytick={-2,-1,0,1,2}, ztick={0,0.5,1},
      grid=major, grid style={grisclaro, very thin},
      axis line style={azulprof!70}, every tick/.style={azulprof!70},
      colormap name=campoclaro, point meta min=0, point meta max=1.3,
      view={-38}{26}, clip=false]
    % superficie con malla
    \addplot3[surf, shader=faceted interp, faceted color=white!40!azulclaro, line width=0.25pt,
              samples=30, domain=-2:2, opacity=0.78] {exp(-(x^2+y^2))};
    % curva de nivel por P (f = e^{-0.89}) sobre la superficie y su proyección
    \addplot3[acento, line width=1.1pt, domain=0:360, samples=90, samples y=0]
      ({0.9434*cos(x)},{0.9434*sin(x)},{0.4107});
    \addplot3[acento, densely dashed, line width=0.6pt, domain=0:360, samples=90, samples y=0]
      ({0.9434*cos(x)},{0.9434*sin(x)},{0});
    % P, su proyección y el gradiente (en el plano z=0)
    \addplot3[only marks, mark=*, mark size=2pt, azulprof] coordinates {(0.8,0.5,0.4107)};
    \addplot3[densely dotted, grispizarra, line width=0.6pt] coordinates {(0.8,0.5,0) (0.8,0.5,0.4107)};
    \draw[-{Stealth[length=6pt]}, line width=1.6pt, acento] (axis cs:0.8,0.5,0) -- (axis cs:0.376,0.235,0);
    % vector tangente a la superficie, de máximo ascenso
    \draw[-{Stealth[length=6pt]}, line width=1.6pt, azulprof] (axis cs:0.8,0.5,0.4107) -- (axis cs:0.376,0.235,0.7977);
    \node[font=\small, azulprof, anchor=east, fill=white, fill opacity=0.85, text opacity=1, inner sep=1.5pt] at (axis cs:-0.1,-0.05,1.0) {máximo ascenso};
    \node[font=\small, azulprof, anchor=south west, fill=white, fill opacity=0.85, text opacity=1, inner sep=1pt] at (axis cs:0.85,0.5,0.45) {$P$};
    \node[font=\small, acento, anchor=north east, fill=white, fill opacity=0.85, text opacity=1, inner sep=1.5pt] at (axis cs:0.3,0.1,0.0) {$\nab f(P)$};
    \node[font=\small, acento, anchor=north west, fill=white, fill opacity=0.85, text opacity=1, inner sep=1.5pt] at (axis cs:-0.5,-1.9,0.0) {curva de nivel $f=c$};
  \end{axis}
\end{tikzpicture}
\caption{Interpretación del gradiente sobre la superficie $z=f(x,y)$. En el
punto $P$, la curva de nivel (en ocre) y el gradiente $\nab f(P)$ (en el plano
del dominio) son perpendiculares. El vector tangente a la superficie con
la misma proyección horizontal apunta en la dirección de máximo ascenso y su
pendiente es $\norm{\nab f(P)}$. El gradiente \emph{no} apunta hacia abajo: en el
plano $xy$ señala hacia donde $f$ aumenta, aquí hacia el origen.}
\label{fig:gradiente-superficie}
\end{figure}


\begin{ejemplo}[Una campana gaussiana]\label{ej:gaussiana}
Sea $f(x,y)=\eul^{-(x^{2}+y^{2})}$, que se parece a una gaussiana. Con la regla
de la cadena $[g(f(x))]'=g'(f(x))f'(x)$, cada derivada parcial produce
un factor $\eul^{-(x^{2}+y^{2})}$ por la derivada del exponente:
\[
  \nab f=\pd{f}{x}\,\ux+\pd{f}{y}\,\uy
  =-2x\,\eul^{-(x^{2}+y^{2})}\,\ux-2y\,\eul^{-(x^{2}+y^{2})}\,\uy
  =-2\,\eul^{-(x^{2}+y^{2})}\,(x\,\ux+y\,\uy).
\]
Es una función con un máximo $f(0,0)=1$ en el origen que decae al alejarnos de
él (\cref{fig:gauss}). Sus curvas de nivel son circunferencias
$x^{2}+y^{2}=-\ln c$, y el gradiente apunta radialmente hacia dentro, hacia el
máximo, perpendicular a esas circunferencias, como exige la
\cref{prop:ortogonal}. Obsérvese además que $\nab f(0,0)=\vc{0}$ (en la cima la
pendiente es nula) y que $\norm{\nab f}\to0$ lejos del origen (la función se
aplana): la pendiente máxima ha de estar en algún punto intermedio, que se
calcula en el \cref{ejer:pendiente}.
\end{ejemplo}

% Figura: f(x,y)=exp(-(x^2+y^2)) y su gradiente (mapa de calor + curvas de
% nivel, superficie 3D, y el campo gradiente en ambas representaciones).
% Los mapas planos se dibujan con anillos concéntricos (f solo depende de r).
\begin{figure}[ht]
\centering
\begin{tikzpicture}[x=1.15cm, y=1.15cm]
  % (a) mapa de calor con curvas de nivel
  \begin{scope}
    \mapaGauss{%
      \foreach \lv/\px/\py in {0.9/0.111/0.305, 0.7/0.561/0.204, 0.5/0.721/-0.416,
                               0.3/0.190/-1.081, 0.2/-0.815/0.972, 0.1/-1.494/-0.264}{%
        \node[font=\tiny, fill=white, fill opacity=0.85, text opacity=1, inner sep=0.6pt]
          at (\px,\py) {\pgfmathprintnumber[fixed,precision=1,use comma]{\lv}};}}
    \node[font=\small\scshape, anchor=south] at (0,2.15) {(a)\; $f(x,y)=\eul^{-(x^{2}+y^{2})}$};
  \end{scope}
  % (c) gradiente sobre el mapa
  \begin{scope}[shift={(0,-6.1)}]
    \mapaGauss{%
      \flechas{-2*exp(-(\x*\x+\y*\y))*\x}{-2*exp(-(\x*\x+\y*\y))*\y}%
        {-2,-1.6,...,2}{-2,-1.6,...,2}{0.75}{0.86}}
    \node[font=\small\scshape, anchor=south] at (0,2.15) {(c)\; campo $\nab f$};
  \end{scope}
\end{tikzpicture}
\hfill
\begin{tikzpicture}
  \pgfplotsset{gsup/.style={superficie, width=0.58\textwidth, height=0.48\textwidth,
      xlabel={$x$}, ylabel={$y$}, zlabel={$f(x,y)$},
      xmin=-2, xmax=2, ymin=-2, ymax=2, zmin=0, zmax=1,
      xtick={-2,-1,0,1,2}, ytick={-2,-1,0,1,2}, ztick={0,0.5,1},
      grid=major, grid style={grisclaro, very thin},
      axis line style={azulprof!70}, every tick/.style={azulprof!70},
      view={-38}{26}, point meta min=0, point meta max=1,
      title style={yshift=3pt}}}
  % (b) superficie con malla
  \begin{axis}[gsup, name=B, title={\small\scshape (b)\; superficie $z=f(x,y)$}]
    \addplot3[surf, shader=faceted interp, faceted color=white!85!azulclaro, line width=0.3pt,
              samples=22, domain=-2:2, opacity=0.97] {exp(-(x^2+y^2))};
  \end{axis}
  % (d) gradiente sobre la superficie
  \begin{axis}[gsup, at={(B.south)}, anchor=north, yshift=-2.4cm,
      title={\small\scshape (d)\; $\nab f$ sobre la superficie}]
    \addplot3[surf, shader=faceted interp, faceted color=white!85!azulclaro, line width=0.3pt,
              samples=22, domain=-2:2, opacity=0.55] {exp(-(x^2+y^2))};
    \addplot3[quiver={u={-1.1*x*exp(-(x^2+y^2))}, v={-1.1*y*exp(-(x^2+y^2))},
                      w={1.1*2*(x^2+y^2)*exp(-2*(x^2+y^2))}, colored=mapped color},
              -{Stealth[length=3pt]}, line width=0.6pt, point meta={2*sqrt(x^2+y^2)*exp(-(x^2+y^2))},
              domain=-1.8:1.8, samples=8, samples y=8] {exp(-(x^2+y^2))+0.015};
  \end{axis}
\end{tikzpicture}
\caption{Campo escalar $f(x,y)=\eul^{-(x^{2}+y^{2})}$ y su gradiente
$\nab f=-2\,\eul^{-(x^{2}+y^{2})}(x\,\ux+y\,\uy)$. (a)~Mapa de calor con
curvas de nivel $f=0{,}1,\,0{,}2,\ldots,0{,}9$ (equipotenciales).
(b)~La misma función vista como superficie $z=f(x,y)$.
(c)~El gradiente apunta hacia el máximo, es ortogonal a las curvas de nivel y
su módulo se anula en el origen y decae a grandes distancias; el color
codifica $\norm{\nab f}$ (de índigo a ámbar). (d)~Sobre la superficie, los
vectores son tangentes a ella y señalan la dirección de máxima pendiente.}
\label{fig:gauss}
\end{figure}


% -------------------------------------------------------------------------
\subsection{Campos vectoriales y campos gradiente}\label{sec:vectoriales}
% -------------------------------------------------------------------------
\begin{definicion}[Campo vectorial]\label{def:vectorial}
Un \emph{campo vectorial} en una región $W\subseteq\R^{n}$ es una función
$\vc{F}:W\to\R^{n}$ que asigna a cada punto un vector,
\[
  \vc{F}(\vc{r})=F_x(\vc{r})\,\ux+F_y(\vc{r})\,\uy+F_z(\vc{r})\,\uz ,
\]
donde las componentes $F_x,F_y,F_z$ son campos escalares.
\end{definicion}

Se visualizan dibujando en cada punto una flecha cuya longitud (y, a menudo, su
color) es proporcional al módulo del campo; esa es la convención que usamos en
todas las figuras de este documento, con una escala cromática de índigo
(módulo pequeño) a ámbar (módulo grande). El campo de velocidades de un fluido,
el campo eléctrico y el gravitatorio son ejemplos físicos. El gradiente de un
campo escalar es siempre un campo vectorial, pero no todo campo vectorial es un
gradiente, lo que motiva la siguiente definición.

\begin{definicion}[Campo gradiente y potencial]\label{def:campo-gradiente}
Un campo vectorial $\vc{F}:W\to\R^{n}$ es un \emph{campo gradiente}
(o \emph{conservativo}) si existe un campo escalar $f:W\to\R$ tal que
$\vc{F}(\vc{r})=\nab f(\vc{r})$ para todo $\vc{r}\in W$. En ese caso $f$ es un
\emph{potencial} (escalar) de $\vc{F}$ en $W$.
\end{definicion}

Los potenciales generalizan la noción de antiderivada de una variable. Hay una
diferencia esencial: toda función continua de una variable en un intervalo
tiene antiderivada, pero para $n\ge2$ existen campos vectoriales continuos
\emph{sin} potencial. Ahora bien, cuando existe, el potencial es esencialmente
único.

\begin{proposicion}[Unicidad del potencial salvo constante]\label{prop:unicidad-potencial}
Si $\vc{F}=\nab f=\nab g$ en una región conexa $W$, entonces $g=f+k$ con $k$
constante.
\end{proposicion}
\begin{proof}
En un dominio conexo, una función $h$ de clase $C^{1}$ tiene gradiente nulo si y
solo si es constante. Como $\nab(f-g)=\vc{0}$, la diferencia $f-g$ es
constante en $W$.
\end{proof}
